\bookchapter{الملحق A: عملية فتح صندوق الهاتف الجديد “خطوة بخطوة”}

\booksection{فحص الصندوق “الثلاثي”:}

\begin{itemize}
  \item \textbf{:سجل دخولك إلى Apple ID أو حساب هواوي أو حساب شاومي وما إلى ذلك.}الإعدادات الأساسية
  \item \textbf{:فعّل "العثور على هاتفي"، واضع رمز قفل شاشة موثوق.}إعداد eSIM
  \item \textbf{:إذا كان هاتفك يدعم eSIM (مثل iPhone وبعض هواتف Android الرائدة)، يمكنك إضافة باقة eSIM في الإعدادات دون إدخال بطاقة SIM فعلية. eSIM مفيد جدًا عند السفر للخارج، حيث يمكنك شراء باقة eSIM من مشغل محلي مباشرة.}نقل البيانات
\end{itemize}

\booksection{التنشيط والإعدادات الأولية “تذكير مهم”:}

\begin{itemize}
  \item \textbf{3. تثبيت التطبيقات "الأمان أولاً":}المتجر الرسمي أولاً
  \item \textbf{:Apple تستخدم App Store، وAndroid يستخدم متجر التطبيقات المدمج في الهاتف.}لا تعطي الصلاحيات عشوائيًا
  \item \textbf{:عند ظهور طلب صلاحية أثناء التثبيت، فكر جيدًا "هل يحتاج هذا التطبيق حقًا إليها؟"}الملحق B: مرجع تمارين عملية
  \item \textbf{تمارين المستوى الأساسي}:نظف تخزين هاتفك، الهدف هو العثور على وحذف مقاطع الفيديو القديمة التي "نسيت ما هي" ولقطات الشاشة المكررة والتطبيقات غير المستخدمة منذ فترة طويلة.
  \item \textbf{تمارين البرمجيات}:اضبط تنظيف تخزين WeChat، المسار عادة [أنا]-[إعدادات]-[عام]-[مساحة التخزين].
\end{itemize}

\booksection{تثبيت التطبيقات “الأمان أولاً”:}

\begin{itemize}
  \item \textbf{:تحقق من صحة البطارية، في iPhone ضمن [إعدادات]-[بطارية]-[صحة البطارية].}نصائح عامة
  \item \textbf{:عند مواجهة أي مشكلة، أعد التشغيل أولاً! هذا المفتاح السحري الذي يحل أكثر من نصف أعطال الهاتف الصغيرة.}الملحق C: جدول مصطلحات تقنية "مبسطة"
\end{itemize}
